% Document control and executive abstract
\section{Document Control \& Executive Abstract}

\subsection{Revision History}
\begin{table}[!htbp]
    \centering
    \small
    \begin{tabular}{llllp{6cm}}
        \toprule
        \textbf{Version} & \textbf{Date} & \textbf{Author} & \textbf{Approved By} & \textbf{Summary} \\
        \midrule
        0.1 & 20 October 2026 & TBD & TBD & Initial modular layout based on the repository artifacts. \\
        \bottomrule
    \end{tabular}
    \caption{Document revision history}
\end{table}

\subsection{Document Purpose}
This High-Level Design Document (HLDD) records the NaviBot system architecture, design decisions, constraints, and planned verification. It is an initial draft based on the files currently available in the repository; TBD items require project-owner input or engineering validation.

\subsection{Document Target Audience}
The intended readers and review framework have not yet been specified. Complete this section when the project audience and submission context are confirmed.

\subsection{Executive Abstract}
NaviBot is planned around a Raspberry Pi 4B, serial YDLIDAR X2, Pi Camera V2 (IMX219), and two classic Arduino Nano motor-controller nodes, one for the front pair and one for the rear pair. The project owner has confirmed a shared USB-RS-485 bus between the Pi and both Nanos. ROS 2 Lyrical Luth on Ubuntu 26.04 arm64 is the proposed Pi software platform. Section 6 describes the camera and LiDAR bring-up, ROS motor-bridge interfaces, control data flow, and firmware contract. Camera streaming and bounded LiDAR scan streaming have been observed on the Pi; sensor calibration and mounting transforms remain open. The current architecture diagram predates this update and still shows one Arduino and a direct UART link. Firmware source, exact components, navigation behavior, and physical bus/motor validation remain to be completed. The project purpose and the network-discovery utility's relationship to the robot also remain to be confirmed.
